\section{Method}\label{sec:method}

This section defines the three quantities the workshop variant uses: the progressive conditional surprise curve $a_k$, the coherence term $C$, and the scalar score $D_{Ca_n} = C \times a_n$.
The full information-theoretic motivation, the alternative excess-entropy summary $E$, the coherence-heterogeneity diagnostic $\sigma_\ell$, and the cross-mode learning analysis appear in Appendices~\ref{app:k0-derivation}--\ref{app:qwen3-comparison}.

\subsection{Progressive Conditional Surprise}\label{sec:progressive-conditioning}

Given prompt $p$ and $n$ responses $r_1, \ldots, r_n$ from policy $\pi$, define
\begin{equation}\label{eq:ak}
    a_k = -\log_2 \theta(r_k \mid r_{<k}, p) \quad \text{for } k = 1, \ldots, n,
\end{equation}
the total surprise of the $k$-th response under base model $\theta$ given the previous $k{-}1$ responses.
The sequence $(a_1, \ldots, a_n)$ is the \textbf{progressive conditional surprise curve}.
We normalize each $a_k$ by the byte count $\|r_k\|$ to get a per-byte (bits/byte) curve that is independent of $\theta$'s tokenizer.
All quantities below are per-byte unless otherwise stated.

The intuition is that as $\theta$ sees more responses from $\pi$, the conditional surprise drops if the responses share patterns ($\theta$'s in-context learning picks up on modes, stylistic regularities, topic patterns) and stays flat if they do not.
For a policy with rich coherent diversity the curve declines and levels off at a positive floor; for a policy producing one repeated mode the curve drops sharply to near zero; for pure noise the curve stays roughly constant at a high value.

In practice we average $a_k$ over uniformly random permutations of the response ordering (so each position averages over all responses and the curve reflects only how $\theta$'s predictions improve with more context, not which response happened to land in slot $k$).
The asymptotic floor is taken to be the last observed point: $a_\infty \approx a_n$.
No fitting or extrapolation step is involved.

\subsection{The Coherence Term}\label{sec:coherence}

The curve alone cannot distinguish a policy producing rich coherent diversity from one producing several individually incoherent modes.
For both, $\theta$ may learn the mode structure from context (so $a_k$ drops meaningfully) and end up at a similar floor.
The distinguishing signal lies in how plausible $\theta$ finds each individual response, independent of the others.

Let $h_\theta(r \mid p) = -\frac{1}{\|r\|}\log_2 \theta(r \mid p)$ be the per-byte cross-entropy of response $r$ under $\theta$ conditioned on the prompt alone.
Define \textbf{coherence} as the geometric mean of the per-byte probabilities:
\begin{equation}\label{eq:coherence}
    C = 2^{-\frac{1}{n}\sum_{i=1}^{n} h_\theta(r_i \mid p)} = \frac{1}{\mathrm{PPL}_\theta(\pi, p)},
\end{equation}
the reciprocal of the geometric-mean per-byte perplexity that $\theta$ assigns to the responses individually.\footnote{For scale, GPT-3 davinci on The Pile \citep{gao2020pile} scored $\sim$0.72 bits/byte and GPT-2 XL $\sim$1.05 bits/byte, giving $C \approx 0.61$ and $C \approx 0.49$ respectively; stronger base models reach lower bits/byte and correspondingly higher $C$.}
Geometric averaging (rather than arithmetic) is what gives $C$ the desired suppression: a single fluent response cannot rescue a set in which the rest are incoherent, because $C$ is driven toward zero by any sample with high per-byte cross-entropy.

\subsection{The Diversity Score}\label{sec:scalar}

The scalar score is the product:
\begin{equation}\label{eq:D-Can}
    D_{Ca_n} = C \times a_n \quad \text{(bits/byte).}
\end{equation}
It can be read as ``how many bits of surprise per byte remain after $\theta$ has learned everything it can from the other responses, weighted by output plausibility.''
The intended edge-case behaviors are summarized in Table~\ref{tab:edge-cases}; empirical metric values on synthetic instances of each scenario appear in Appendix~\ref{sec:scenario-validation} (Table~\ref{tab:scenarios}).

\begin{table}[t]
\centering
\small
\caption{Intended edge-case behavior of $D_{Ca_n} = C \times a_n$ on the five appendix scenarios (Appendix~\ref{sec:scenario-validation}).
The product correctly suppresses pure noise and incoherent multi-mode (via $C$) and correctly suppresses one-mode sets (via $a_n$); multi-mode coherent is the intended winner.
Empirical numbers under GPT-2 and Qwen2.5-3B appear in Table~\ref{tab:scenarios}.}
\label{tab:edge-cases}
\begin{tabular}{@{}lccc@{}}
\toprule
Scenario              & $C$  & $a_n$ & $D_{Ca_n}$       \\
\midrule
Pure noise            & low  & high  & low (via $C$)    \\
Multi-mode incoherent & low  & high  & low (via $C$)    \\
Multi-mode coherent   & high & high  & \textbf{high}    \\
One-mode              & high & low   & low (via $a_n$)  \\
Mixed                 & mid  & high  & mid              \\
\bottomrule
\end{tabular}
\end{table}

The score depends on $p$, $\pi$'s responses, and $\theta$.
Diversity is therefore a property of (responses, prompt, scoring model), not of the policy in isolation: the same response set can score differently under a stronger or weaker $\theta$.
This is by design: $\theta$'s in-context learning capability is the lens through which diversity is measured, and the score sharpens automatically as base models improve.
When comparing policies across a prompt suite, $D_{Ca_n}$ can be averaged over prompts or differenced; see Appendix~\ref{app:aggregation}.

